\section{问题分析}
\label{sec:2}

\subsection{总体思路}
\label{sec:2.1}

四问按「测量数据—描述损失—优化训练—解释能力」层层递进：问题一从文档质量 $Q(x)$ 聚合到领域质量 $Q_i$，再按配比得到语料质量 $Q(\bp)$；问题二在 $L_0(N,D)$ 上引入质量与配比，经典律正是广义律在 $Q=1,\bp=\bp_{\rm ref}$ 时的切片；问题三以标度律为目标，加入预算、提质与注意力成本约束求最优配置；问题四引入时间与 Benchmark 度量，把损失变化转化为能力前沿的解释与条件预测。各问的冻结结果通过 P1–P4 接口传递，未通过留出检验的配比通道不进入下游主优化。四问的参数、约束与接口传递关系在 \ref{sec:5.1} 的接口约定与 \ref{sec:5.2}–\ref{sec:5.5} 各问的模型建立中逐条说明。

\subsection{各问题的关键难点}
\label{sec:2.2}

\textbf{问题一：从异质指标得到唯一语料质量。} logits、规则量与 DSIR 的方向和尺度各异，长度与领域基线会干扰质量判断，17 域中又有 11 域缺少直接评分。为避免冗余信号重复计权、避免单项极端指标支配整篇分数，本文采用 A1 冻结的域内预处理、CRITIC 权重与 Huber 中心；冲突在评分之后按来源诊断，不反向改分。领域质量以映射加文本桥接补齐，配比因份额闭合约束采用中心对数比建模，并以 A6–A11 检验跨尺度边界。

\textbf{问题二：统一跨附件质量口径并确定可推广的函数形式。} A、B 两组损失的分词器与验证集可能不同，绝对值不可直接对齐，因此质量项不改动 B1 骨架的不可约项 $E$，只乘在可约部分上，另加一项在 $Q=1$ 处为零的地板。质量乘子的形状（幂次、指数、线性）不事先指定，本文在同一骨架上比较 M0–M4 与有无质量地板（不随规模消失的加项 $k_0(1-Q)$）的嵌套族，用 B6 训练、B7 新增点留出。主式指定为幂次，以便把质量解释为有效规模；交叉验证表现更优的线性对照一并报告。配比效应仅有 1M 数据，无法识别其随规模的斜率，因此不估计尺度项，只设两个零新增参数的强度假设，再用 A6–A11 的观测斜率作事后比较。

\textbf{问题三：处理非线性约束与边界转移。} 三项成本都与 $D$ 成正比，而损失对 $D$ 严格递减，故预算取等、可解析消去 $D$，把问题降为 $(\log_{10}N,Q)$ 的二维有界优化。质量约束是否活跃构成质量状态的边界分类，其变化即最优策略的质变。三类提质成本函数的水平与边际曲率各不相同，单个成本项为凸不保证联合优化为凸，因此保留独立粗网格核对，并区分局部一阶根与全局状态切换。

\textbf{问题四：界定能力演进的解释边界。} Loss 与 Benchmark 得分尺度不同、任务会饱和、规模与时间共变、C3 口径混杂。本文采用可检验的 S 形损失—得分桥接与连续时间项，在一致的开放权重与模型类型口径下分解增长，并做窗口敏感性与回测。时间项只表示控制规模状态后的剩余关联，不直接解释为已识别的因果技术进步。

\subsection{数据使用与证据分层}
\label{sec:2.3}

各附件的角色与证据性质见表~\ref{tab:attach}，表中角色以对应 Spec 为准。

\begin{table}[H]
\centering\footnotesize
\caption{附件角色与证据性质（引自各问 Spec，非数据产物）}
\label{tab:attach}
\begin{tabularx}{\textwidth}{@{}p{1.25cm}X p{3.8cm}p{3.05cm}@{}}
\toprule
附件 & 内容 & 角色 & 证据性质 \\
\midrule
A1 & 七域质量信号抽样 51,230 篇 & 拟合与标准化参照（唯一） & 真实 \\
A2 / A3 & arxiv 17,523 / github 203,752 篇 & 冻结评分与诊断的扩展检验 & 真实 \\
A4 / A5 & 1M 配比与 13 域 Loss 各 512 组 & 配比主模型训练（唯一） & 真实 \\
A6 / A7 & 1M 配比与 Loss 各 256 组 & 同尺度留出 & 真实 \\
A8 / A9 & 60M 配比与 Loss 各 256 组 & 跨尺度直接检验 + 尺度修正标定 & 真实 \\
A10 / A11 & 1B 配比与 Loss 各 64 组 & 独立检验，不参与任何估计 & 真实 \\
A12–A15 & 10B/70B 配比（各 63 组）与 Loss & 估算/外推情景对照 & 子集 / 外推，非观测 \\
A16 & 17 域$\to$7 质量域映射 & 映射参照 & 人工整理 \\
A17 & 域统计摘要 & 背景与辅助 & 真实，未接入求解 \\
A18 & 17 域原始文本 138,034 条（本轮有效 31,340 篇用于桥接预测） & 桥接预测输入 & 真实 \\
B1 & Pythia 训练轨迹 1,176 行 & 经典基线拟合 & 真实训练设定，损失按公开公式生成（见审计） \\
B2 & Cerebras 轨迹 1,029 行 & 族外趋势对照 & \textbf{半合成} \\
B3 & 插值轨迹 4,000 行 & 轨迹对照 & 插值（B1 派生） \\
B4 / B5 & 跨族基准 57 行 / 文献点 44 行 & 外部趋势对照（$Q=1$） & 真实 / 文献整理 \\
B6 & NQ 质量网格 360 行 & 质量项参数估计（唯一） & \textbf{半合成} \\
B7 & 扩展网格 450 行（新增 90 点） & 留出检验 & \textbf{半合成}，$Q$ 方向插值 \\
B8 & 大规模矩阵 1,704 行 & 方向冲突压力检查 & \textbf{半合成}，与 B6/B7 方向相反 \\
B9 / B10 & 100B+ 配置 132 行 / 预估 Loss 128 行 & 外推参照 & 真实元数据 / \textbf{估算} \\
B11 / B12 & 模型族元数据 / Pythia 检查点索引 & 来源审计 & 真实 \\
C1 / C2 & 榜单 4,576 行 + Epoch 匹配列 & 能力与发布日期 & 真实（同批榜单，非两套样本） \\
C3 & 历史时序 4,599 行 & 降权敏感性 & \textbf{混合口径} \\
C4 & 全量模型元数据 3,523 行 & 规模、算力、开放权重 & 真实 \\
C5 / C6 & 桥接表 43 / 75 行 & 损失—得分映射 & \textbf{混合口径}，分层使用 \\
C7 & 架构元数据 45 行 & 上下文长度档位 & 真实 \\
C8 & 逐任务 JSON 1,958 个 & 逐任务聚合与饱和度 & 真实 \\
\bottomrule
\end{tabularx}
\end{table}
